\section{Related Work}
\label{sec:related}

\para{Multi-agent LLM debate.}
\citet{du2023debate} establish the foundational ``society of minds'' paradigm,
where multiple instances of the same LLM propose, debate, and update answers over
multiple rounds. They show that homogeneous multi-agent debate improves mathematical
reasoning on GSM8K by 11\% over single models and scales with agent count.
Critically, they identify that RLHF-trained agents are ``agreeable''---they readily
accept critiques even when originally correct---a finding directly relevant to
our debate condition.
\citet{liang2023mad} identify the ``Degeneration-of-Thought'' problem: once an LLM
establishes high confidence, single-model self-reflection fails to generate novel
reasoning. Their Multi-Agent Debate (MAD) framework uses two adversarial agents
managed by a judge and improves counter-intuitive arithmetic reasoning.
Building on these works, we study debate between agents of \emph{different capability levels},
rather than multiple instances of the same model.

\para{Heterogeneous model ensembles.}
\citet{chen2023reconcile} introduce \textsc{ReConcile}, the first round-table discussion
framework using genuinely heterogeneous models (GPT-3.5, PaLM-2, Claude).
They demonstrate that diversity across model families is critical to superior
performance (+11.4\% on CommonsenseQA over best baselines) and introduce
confidence-weighted voting for consensus.
Unlike \citet{chen2023reconcile}, we focus specifically on \emph{within-family}
capability diversity (small vs.\ large variants of the same model), controlling
for provider-level differences to isolate the effect of model scale.
\citet{wu2023heterogeneous} provide the theoretical foundation through
negative correlation learning analysis, showing that heterogeneous deep ensembles
with high diversity boost adversarial robustness in computer vision.
We apply this theoretical lens to LLM ensembles and measure pairwise error
correlation directly.

\para{Scaling and voting in homogeneous ensembles.}
\citet{li2024moreagents} show that LLM performance scales monotonically with
agent count in a sampling-and-voting (Agent Forest) framework,
even for homogeneous ensembles.
The improvement is stronger on harder tasks.
This work provides an important baseline: any claimed benefit of diversity
must exceed what is achievable by simply scaling homogeneous sampling.
In our study, we control for agent count (three agents per ensemble condition)
to isolate diversity effects from scale effects.

\para{Diversity metrics and ensemble optimization.}
\citet{tekin2024topla} introduce the focal diversity metric for quantifying
diversity-performance correlation in LLM ensembles and propose
diversity-optimized pruning to select the top-performing sub-ensemble.
Their \textsc{LLM-TOPLA} framework achieves +2.2\% on MMLU over strong baselines
by selecting diverse sub-ensembles.
We adopt their conceptual framework of measuring pairwise error correlation
as a diversity index but evaluate it in the context of capability-level
rather than cross-architecture diversity.

\para{Challenges to the diversity hypothesis.}
\citet{rosales2025diverse} present the most direct challenge to our hypothesis:
they show that question interpretation diversity (rephrasing the question for a
single model) often outperforms model diversity across GPT and Llama families
on BoolQ, StrategyQA, and PubMedQA.
Their key insight---that model diversity benefits only when models have
\emph{complementary failure modes}, not just independent errors---motivates
our investigation into whether within-family models share failure modes
due to shared safety alignment.

\para{Controlled studies of debate factors.}
\citet{wu2025debate} conduct a controlled study of multi-agent debate using
Knight-Knave-Spy logical puzzles with verifiable ground truth.
They find that group diversity and intrinsic reasoning strength dominate
debate outcomes, while structural parameters (order, confidence visibility)
offer minimal gains.
Diverse teams overturn incorrect consensus while homogeneous teams create
echo chambers.
This finding motivates our investigation of whether capability diversity
(small vs.\ large within the same family) translates to the kind of beneficial
debate diversity observed with cross-family models.

\para{Positioning.}
Unlike prior work, we specifically target the \emph{within-family capability diversity}
question with a custom adversarial benchmark designed to expose shared failure modes,
and we measure both accuracy and pairwise error correlation.
This provides complementary evidence to the cross-family results of
\citet{chen2023reconcile} and the cross-model debate results of \citet{wu2025debate}.
